\subsection{CT4 --- Charging schemes}\label{ct:4}\label{charging-schemes}
\ctsignature{\(P_{\ast\to4}\to \ctcert{C4};\ \ctint{P_{4\to9},P_{4\to14}};\ \ctrec{P_{4\to13}}\).}

CT4 is an exact sequential machine for a deterministic demand--payer ledger.
Its \texttt{Framework} fixes dependent demand and payer types, eligibility,
canonicality, fibre capacity, the global demand--capacity comparison, the C4
claim, and downstream datum types.  A validated \texttt{Input} contains
\(\langle G,\mathfrak B,\pi,C,\ell\rangle\), where \(\pi\) is an explicit
partial assignment function.  Scope readiness, canonicality, totality, bounded
fibres, and the final comparison are distinct contracts.

The core plan has five independent fields:
\texttt{scope}, \texttt{assignment}, \texttt{availability},
\texttt{fibres}, and \texttt{comparison}.  The assignment node is a
single-successor certification boundary.  It constructs
\texttt{AssignmentState}, whose \texttt{AssignmentCertificate} proves that
the supplied function implements the declared deterministic tie-breaking
rule.  A missing proof of canonicality rejects the plan before execution; it
does not create an untyped recovery edge.

The availability node consumes the exact assignment state.  It either builds a
\texttt{CT13Payload} containing a concrete unassigned demand and proof that no
eligible payer exists, or constructs \texttt{TotalAssignmentState}, which
proves that every demand is assigned to an eligible payer.  The fibre node is
indexed by this totality proof and returns either an overload payload for CT9
or \texttt{BoundedFibreState}.  The comparison node is indexed by the bounded
state and returns a C4 certificate or a CT14 aggregate residual.  Thus
determinism, payer availability, no-overcount, and the numerical comparison
can be implemented and verified independently.

\texttt{Port} and \texttt{HandoffPlan} contain only consumer-side acceptance
proofs for CT13, CT9, and CT14.  They are applied after \texttt{runCore} and
cannot alter the mathematical result.  Theorems \texttt{run\_verified},
\texttt{run\_total}, and \texttt{run\_trace\_valid} establish soundness,
relative totality, and exact trace validity.

The soundness contracts are assigned as follows.
\begin{itemize}
\tightlist
\item \textbf{S-Def} fixes the demand family, payer family, assignment, capacity datum, and lower-bound datum.
\item \textbf{S-Det} is discharged only by \texttt{AssignmentCertificate}.
\item \textbf{S-Dich} is discharged by the scope, availability, fibre, and comparison constructors.
\item \textbf{S-Comp} is carried by \texttt{BoundedFibreState} and \texttt{C4Certificate}.
\item \textbf{S-Rout} and \textbf{S-Trig} are discharged by the indexed CT13, CT9, and CT14 payloads and the separate handoff plan.
\end{itemize}

\begin{figure}[p]
\centering
\vspace*{-1.70cm}
\resizebox{0.94\textwidth}{!}{%
\begin{tikzpicture}[x=1cm,y=1cm,every node/.style={font=\scriptsize}]
\node[bcwide] (entry) at (0,0) {{\tiny\ttfamily CT4.entry}\\\textbf{Validated ledger}\\\((G,\mathfrak B,\pi,C,\ell)\)};
\node[bcdec,bclemma,text width=3.35cm,minimum width=4.10cm] (scope) at (0,-2.35) {{\tiny\ttfamily CT4.decide.scope}\\\textbf{Demand family admitted?}};
\node[bcscopefail,text width=3.20cm,minimum width=3.20cm] (scopeout) at (7.25,-2.35) {{\tiny\ttfamily CT4.terminal.scope}\\\textbf{Scope exit}};
\node[bcwide,bclemma] (assignment) at (0,-4.80) {{\tiny\ttfamily CT4.certify.assignment}\\\textbf{Canonical assignment certificate}\\explicit deterministic tie-breaking};
\node[bcdec,bclemma,text width=3.35cm,minimum width=4.10cm] (availability) at (0,-7.30) {{\tiny\ttfamily CT4.decide.availability}\\\textbf{Every demand has a payer?}};
\node[bcroute,bcrouteneutral,bcdesign,text width=3.20cm,minimum width=3.20cm] (ct13) at (7.25,-7.30) {{\tiny\ttfamily CT4.terminal.ct13}\\\textbf{CT 13 exit}\\missing payer};
\node[bcdec,bclemma,text width=3.35cm,minimum width=4.10cm] (fibres) at (0,-10.15) {{\tiny\ttfamily CT4.decide.fibres}\\\textbf{All fibres bounded?}};
\node[bcroute,bcrouteneutral,bcoutputroute,text width=3.20cm,minimum width=3.20cm] (ct9) at (-6.70,-10.15) {{\tiny\ttfamily CT4.terminal.ct9}\\\textbf{CT 9 exit}\\overload fibre};
\node[bcdec,bctheorem,text width=3.35cm,minimum width=4.10cm] (comparison) at (0,-13.05) {{\tiny\ttfamily CT4.decide.comparison}\\\textbf{Demand exceeds capacity?}};
\node[bcclosed,text width=3.00cm,minimum width=3.00cm] (c4) at (-4.20,-16.05) {{\tiny\ttfamily CT4.terminal.c4}\\\textbf{C4}};
\node[bcroute,bcrouteneutral,bcoutputroute,text width=3.30cm,minimum width=3.30cm] (ct14) at (4.20,-16.05) {{\tiny\ttfamily CT4.terminal.ct14}\\\textbf{CT 14 exit}\\aggregate residual};
\draw[bcarr] (entry)--(scope);
\draw[bcscopearr] (scope)--node[bclabel,above]{exit}(scopeout);
\draw[bcarr] (scope)--node[bclabel,right]{ready}(assignment);
\draw[bcarr] (assignment)--node[bclabel,right]{canonical}(availability);
\draw[bcrecoverarr] (availability)--node[bclabel,above]{\(P_{4\to13}\)}(ct13);
\draw[bcarr] (availability)--node[bclabel,right]{total}(fibres);
\draw[bchandoff] (fibres)--node[bclabel,above]{\(P_{4\to9}\)}(ct9);
\draw[bcarr] (fibres)--node[bclabel,right]{bounded}(comparison);
\draw[bcarr] (comparison)--node[bclabel,above,sloped]{exceeds}(c4);
\draw[bchandoff] (comparison)--node[bclabel,above,sloped]{\(P_{4\to14}\)}(ct14);
\end{tikzpicture}%
}
\caption[Closure-tactic 4: exact charging machine]{Closure-tactic 4 as an exact sequential machine.  The assignment certificate, totality state, and bounded-fibre state are separate proof boundaries.  Each later node is indexed by the exact earlier proof, so an overload, missing-payer witness, C4 certificate, or aggregate residual cannot be attached to another ledger.  Monospaced identifiers are shared by Lean, JSON, and executable traces.}
\label{fig:ct4-charging-schemes}
\end{figure}
\FloatBarrier
